\section{Centered moments of the Schur label}
\label{sec:schur-label-moments}

The signed label-moment argument in
\cite[Section~7, lines~227--237]{OpenAI2026AreaLaw} first compares the label
observable with the supported surprisal. The estimates below give this
comparison for an arbitrary real center. The subsequent bounds for the
centered surprisal of independent copies are a separate argument.

Let a finite group act by permutations on a finite-dimensional Hilbert space,
let $\rho\ge0$ commute with that action, and let $P_\lambda$ be the central
projection belonging to the irreducible representation of dimension $d_\lambda$.
Write
\begin{align}
 F&=\sum_\lambda (\log d_\lambda)P_\lambda,\\
 L&=-\log\rho,\\
 M_A(t;s)&=\operatorname{ReTr}\bigl(\rho e^{t(A-sI)}\bigr).
\end{align}
The logarithm of $\rho$ is evaluated on its support and is zero on its kernel.

\begin{theorem}[Positive centered comparison]
\label{thm:schur-label-positive-centered}
\lean{PermutationRepresentation.re_trace_mul_exp_centered_labelEntropy_le}
\uses{thm:schur_label_surprisal}
\leanok
If $\operatorname{Tr}\rho=1$, then for every $s\in\mathbb R$ and $u\ge0$,
\begin{align}
 M_F(u;s)\le M_L(u;s).
\end{align}
\end{theorem}

\begin{proof}
\leanok
Use the joint spectral resolution of $\rho$ and the central projections.
On a nonzero joint eigenspace with positive eigenvalue $r$, invariance and
normalization give $r d_\lambda\le1$, so $\log d_\lambda\le-\log r$.
The desired comparison follows by monotonicity of the scalar exponential.
The terms with $r=0$ have zero trace weight. Thus only the supported
comparison is used; an inequality on the entire ambient space is unnecessary.
\end{proof}

\begin{theorem}[Negative centered comparison]
\label{thm:schur-label-negative-centered}
\lean{PermutationRepresentation.re_trace_mul_exp_neg_centered_labelEntropy_le}
\uses{thm:weighted-trace-holder-binary}
\leanok
Without requiring $\operatorname{Tr}\rho=1$, for all $s,u\in\mathbb R$,
\begin{align}
 M_F(-u;s)\le
 \sqrt{M_L(-2u;s)\operatorname{ReTr}\bigl(\rho e^{2u(L-F)}\bigr)}.
\end{align}
\end{theorem}

\begin{proof}
\leanok
The Hermitian operators $-2u(L-sI)$ and $2u(L-F)$ commute, since both are
functions of the joint spectral resolution. Their average is $-u(F-sI)$.
Apply the weighted trace inequality with weights $1/2,1/2$. The individual
exponential traces are nonnegative, so the product of their square roots
is the square root of their product.
\end{proof}

\begin{theorem}[Copy-permutation moment transfer]
\label{thm:schur-label-copy-centered}
\lean{TensorPower.centered_labelEntropy_moments_le}
\uses{thm:schur-label-positive-centered,thm:schur-label-negative-centered,
thm:schur_label_surprisal}
\leanok
Let $\rho$ be a permutation-invariant density matrix on $Q^{\otimes k}$,
where $\dim Q=q$. For every $s\in\mathbb R$ and $0\le u\le1/2$,
\begin{align}
 M_F(u;s)&\le M_L(u;s),\\
 M_F(-u;s)&\le\sqrt{(k+1)^{q^2}M_L(-2u;s)}.
\end{align}
\end{theorem}

\begin{proof}
\leanok
Apply the preceding comparisons and the Schur remainder estimate
$\operatorname{ReTr}(\rho e^{b(L-F)})\le(k+1)^{q^2}$ with $b=2u\le1$.
The other factor is nonnegative, so multiplication and the square root
preserve the bound.
\end{proof}
